\subsection{Momentum dependence of the pole residues}
\label{sec:lovelock-residues}

Use the retarded normalization \(\mathcal G=\mathcal V/\mathcal S\),
where \(A_{\rm in}\sim\mathcal S r^{\nu-1}+\mathcal V r^{-\nu-1}\)
in the two boundary branches. First take \(\nu\) noninteger. At
zero momentum the horizon solution and its source coefficient are
\begin{align}
 H_\omega(x)&=x^{-{\ii w}/2}(1-x)^h
       {}_2F_1(\mathsf a,\mathsf b;\mathsf c;x),\notag\\
 \mathsf a&=h-s_+,\qquad \mathsf b=h+s_-,
 \qquad \mathsf c=1-\ii w,\notag\\
 \mathcal S(\omega)&=2^{1-h}
  \frac{\Gamma(\mathsf c)\Gamma(\nu)}
       {\Gamma(\mathsf a)\Gamma(\mathsf b)} ,
 \label{eq:lovelock-horizon-general}
\end{align}
where \(x=(r^2-1)/(r^2+1)\). Here \(\mathcal S\) refers to a
unit leading horizon coefficient in \(x\).

Let a dot denote the derivative with respect to \(k^2\) at
fixed \(\omega\) and \(k=0\). The perturbation of the ingoing
solution with unit source satisfies
\(\mathcal L_\omega\dot{\widehat A}=\widehat A/r\), with
\(\widehat A=H_\omega/\mathcal S\) and \(\mathcal L_\omega\)
the radial operator in \eqref{eq:lovelock-k-radial} at \(k=0\).
The Wronskian at the boundary is \(-2\nu\dot{\mathcal G}\).
The horizon term vanishes for \(\operatorname{Im}\omega>0\),
where the ingoing solution decays. For \(0<\nu<1\) the
radial integral converges at the boundary. The resulting response
derivative is
\begin{align}
 \dot{\mathcal G}(\omega)
 &=-\frac{\mathcal I(\omega)}{2\nu\mathcal S(\omega)^2},\notag\\
 \mathcal I(\omega)&=
 \operatorname{FP}\!\int_0^1
 \frac{x^{-\ii w}(1-x)^\nu}{1+x}
 {}_2F_1(\mathsf a,\mathsf b;\mathsf c;x)^2\,\dd x .
 \label{eq:lovelock-response-derivative}
\end{align}
The finite part continues the horizon integral to the pole
region and the boundary integral in \(\nu\). For noninteger
\(\nu>1\) it removes the divergent powers of the cutoff
from the source branch. At integer \(\nu\) there are local
logarithmic subtractions. Polynomial contact terms have no
effect on the pole coefficients.

For \(N=n+h\), the zero-momentum residue on the positive-real-part
branch follows directly from \eqref{eq:detuned-g}:
\begin{equation}
 \mathcal R_n^+=-\frac{4\,2^\nu}{(1+\ii)\Gamma(\nu)}
 \frac{(\nu+1)_n}{n!}
 \frac{\Gamma(h-\ii N)}{\Gamma(1-h-\ii N)} .
 \label{eq:lovelock-exact-residue}
\end{equation}
At \(\omega_n^0=2(1-\ii)N\), \(\mathsf a=-n\), and the factor
\(1/\Gamma(-n-\nu)\) cancels \(\Gamma(-\nu)\) in the
residue. This gives the finite expression
\eqref{eq:lovelock-exact-residue} also when \(\nu\) is integer.
The negative-real-part residue is
\(\mathcal R_n^-=-[\mathcal R_n^+]^*\).

The expansion of \eqref{eq:lovelock-response-derivative} near
\(\omega_n^0\) contains a double and a simple pole,
\begin{equation}
 \dot{\mathcal G}(\omega)
 =\frac{\mathcal R_n^+ b_n}{(\omega-\omega_n^0)^2}
    +\frac{\dot{\mathcal R}_n^+}{\omega-\omega_n^0}+O(1),
 \label{eq:lovelock-derivative-principal-part}
\end{equation}
where \(\dot{\mathcal R}_n^+\) is the change of the residue
with the pole moving. The two coefficients are
\begin{align}
 \mathcal R_n^+ b_n
 &=-\left.\frac{\mathcal I}{2\nu(\mathcal S')^2}
                \right|_{\omega_n^0},\notag\\
 \dot{\mathcal R}_n^+
 &=-\left.\frac{\mathcal I'-\mathcal I\mathcal S''/\mathcal S'}
                   {2\nu(\mathcal S')^2}\right|_{\omega_n^0}.
 \label{eq:lovelock-residue-derivatives}
\end{align}
Primes denote \(\omega\) derivatives. Thus the continued integral
fixes both pieces of the response at order \(k^2\).
Expanding the shifted pole gives precisely the double-pole
coefficient determined by \eqref{eq:lovelock-exact-dispersion}.

\begin{table}[t]
 \centering\small
 \begin{tabular}{ccccc}
 \hline
 \(\nu\) & \(n\) & \(\dot{\mathcal R}_n^+\) &
 Double-pole error & Direct response error \\
 \hline
 0.5 & 0 & $0.0859270+0.0975837\ii$ & 7.98e-15 & 1.30e-11 \\
0.75 & 2 & $-0.0662687+0.2797082\ii$ & 9.33e-14 & 1.46e-12 \\
1.25 & 3 & $1.5665266+4.1621940\ii$ & 7.68e-10 & 2.13e-09 \\
\hline

 \end{tabular}
 \caption{Momentum derivatives of the retarded residues. The
 double-pole error compares the integral with \(\mathcal R_n^+ b_n\).
 The final column gives the maximum absolute error for the
 simple and double coefficients from the direct finite-momentum response.}
 \label{tab:lovelock-residues}
\end{table}

In table~\ref{tab:lovelock-residues} we evaluate
\eqref{eq:lovelock-response-derivative} using convergent
hypergeometric series on \(0\leq x\leq1/2\) and
\(1/2\leq x\leq1\), integrating each power with its
analytic finite part. We extract the principal part by
quadrature on a circle around the pole. As an independent
check, we solve the full finite-momentum equation for the
source and response on the same circle and extrapolate the
centered difference in \(k^2\).

\subsection{Primitive times from the zero-momentum mode sum}

The residues \eqref{eq:lovelock-exact-residue} give the
retarded kernel for real \(t>0\) in the form
\begin{equation}
 \mathcal G_R(t)=-\ii\sum_{n=0}^{\infty}
 \big[\mathcal R_n^+ e^{-\ii\omega_n^+ t}
      +\mathcal R_n^- e^{-\ii\omega_n^- t}\big],
 \qquad \omega_n^\pm=\pm 2N-2\ii N .
 \label{eq:lovelock-mode-sum}
\end{equation}
One way to close the inverse Fourier contour is to divide
by \((\omega+\ii\Lambda)^J\) for \(\Lambda>0\) and a sufficiently
large integer \(J\). Stirling asymptotics and the gamma reflection
formula bound the ratios in \eqref{eq:detuned-g} by powers of
\(|\omega|\) on arcs between successive pole radii \cite{NIST:DLMF}.
The reflected sine factors are bounded there since the
contours stay a fixed distance from the pole sequences.
The regulated arcs vanish for \(t>0\).
Applying \((\ii\partial_t+\ii\Lambda)^J\) kills the artificial
pole terms and restores \eqref{eq:lovelock-mode-sum}.
Integer \(\nu\) is obtained by the contact-subtracted limit,
whose frequency growth includes at most a logarithm multiplying
the power bound.

The sum converges locally uniformly in
\(\operatorname{Re}t>|\operatorname{Im}t|\).
On its upper edge, at \(t_m=\tau m(1+\ii)\), the
positive-real-part modes have a common phase
\(e^{-2\pi\ii m h}\) and the opposite modes still decay.
In particular, if
\begin{equation}
 C_{\rm R}=-\frac{4\,2^\nu e^{-\ii\pi\nu/2}}
       {(1+\ii)\Gamma(\nu)\Gamma(1+\nu)} ,
\end{equation}
then \(\mathcal R_n^+\sim C_{\rm R}N^{2\nu}\), and the leading
singular term obtained from the mode sum is
\begin{equation}
 \mathcal G_R(t)\sim
 -\ii C_{\rm R} e^{-2\pi\ii m h}\,
 \frac{\Gamma(2\nu+1)}
 {[2(1+\ii)(t-t_m)]^{2\nu+1}} .
 \label{eq:lovelock-retarded-time}
\end{equation}
This follows by rescaling the large-overtone sum, for
example as \(t=t_m+\epsilon\) with \(\epsilon>0\).
The lower edge is obtained by conjugation. The retarded
kernel therefore reproduces the primitive optical times
and their leading powers directly from the pole weights.
In the normalization \(\rho=2\operatorname{Im}\mathcal G\),
the vacuum coefficient is
\(C_\nu=2\pi/[4^\nu\Gamma(\nu)\Gamma(1+\nu)]\).
Dividing the leading term of \eqref{eq:lovelock-retarded-time}
by \(\ii\) gives the Wightman coefficient in
\eqref{eq:primitive-coefficient} at the same time and on the
corresponding approach from the two domains of continuation.

For positive integer \(\nu\), \(\mathcal R_n^\pm\) are
polynomials in \(n\). Their mode sums are meromorphic in \(t\).
For the tuned case \(\nu=1\), set \(q_\pm=e^{-2(1\pm\ii)t}\).
The exact sum reads
\begin{equation}
 \mathcal G_R(t)=-4\ii\left[
 (1+\ii)\frac{q_+(1+q_+)}{(1-q_+)^3}
 -(1-\ii)\frac{q_-(1+q_-)}{(1-q_-)^3}
 \right].
 \label{eq:lovelock-retarded-tuned}
\end{equation}
We check this expression by Fourier transforming the
commutator spectral density \(2\operatorname{Im}\mathcal G\)
using \eqref{eq:lovelock-spectral}. The pole and frequency
sums agree at real \(t>0\). The mixed thermal images are
a feature of the Wightman function in
section~\ref{sec:mixed-time}.
